% Artigo do projeto final de FGA0315 Qualidade de Software 1 (UnB/FGA, 2026/2).
% Modelo SBC (sbc-template.sty e sbc.bst).
% Compilação: pdflatex main && bibtex main && pdflatex main && pdflatex main
% (ou tectonic -X compile main.tex).

\documentclass[12pt]{article}

\usepackage{sbc-template}
\usepackage[T1]{fontenc}
\usepackage[utf8]{inputenc}
\usepackage[brazilian]{babel}
\usepackage{graphicx,url}
\usepackage{multirow}
\usepackage{booktabs}
\usepackage{placeins}
\usepackage{microtype}
\usepackage{array}
\newcolumntype{P}[1]{>{\raggedright\arraybackslash}p{#1}}
% Mantém tabelas e figuras próximas à chamada, sem páginas de floats esparsas.
\renewcommand{\topfraction}{0.9}
\renewcommand{\bottomfraction}{0.8}
\renewcommand{\textfraction}{0.08}
\renewcommand{\floatpagefraction}{0.8}
\setcounter{topnumber}{4}
\setcounter{bottomnumber}{4}
\setcounter{totalnumber}{8}
\usepackage{amsmath}
\usepackage{tikz}
\setlength{\textfloatsep}{6pt plus 2pt minus 2pt}
\setlength{\intextsep}{6pt plus 2pt minus 2pt}
\setlength{\abovecaptionskip}{3pt}
\setlength{\belowcaptionskip}{3pt}
\usetikzlibrary{positioning,fit,backgrounds,arrows.meta,shapes.geometric}

\sloppy

% Permite quebra de linha em URLs longas
\def\UrlBreaks{\do\/\do-\do\_}

% Uma seção obrigatória ausente deve interromper a compilação.
\newcommand{\incluir}[1]{\input{#1}}

\title{Avaliação Comparativa de Qualidade em Serviços Digitais de Governo:\\
  Portal de Serviços $\times$ Aplicativo na Consulta de Infrações do Detran-DF}

\author{Daniel de Oliveira Lira, Heitor Macedo Ricardo, Hugo Freitas Silva,\\
  Pedro Druck Montalvão Reis, Pedro Henrique Ferreira Xavier,\\
  Philipe Amâncio Reis Caetano}

\address{Faculdade de Ciências e Tecnologias em Engenharia (FCTE)\\
  Universidade de Brasília (UnB) -- Brasília, DF -- Brasil}

\begin{document}

\maketitle

\begin{center}
  \small Repositório com dados e evidências:\\
  \url{https://github.com/FCTE-Qualidade-de-Software-1/2026-2_T02_KATHERINE_JOHNSON}
\end{center}

\incluir{secoes/resumo}
\incluir{secoes/introducao}

\section{Referencial Teórico}
\label{sec:referencial-teorico}
\incluir{secoes/referencial-square}
\incluir{secoes/referencial-gqm}

\incluir{secoes/metodologia}

\subsection{Fase 1: Estabelecimento dos Requisitos da Avaliação}
\label{sec:fase1}
\incluir{secoes/fase1-proposito}
\incluir{secoes/fase1-produto}
\incluir{secoes/fase1-priorizacao}

\FloatBarrier
\subsection{Fase 2: Especificação da Avaliação}
\label{sec:fase2-gqm}
\incluir{secoes/fase2-objetivos}
\incluir{secoes/fase2-metricas}

\FloatBarrier
\incluir{secoes/fase3-plano}

\FloatBarrier
\incluir{secoes/fase4-analise}

\incluir{secoes/fase4-conclusoes}

\FloatBarrier
\incluir{secoes/declaracao-ia}

\FloatBarrier
\bibliographystyle{sbc}
\bibliography{referencias}

% Apêndice de contribuição incluído no limite total de páginas
\appendix
\incluir{secoes/contribuicao}

\end{document}
